\documentclass{article}
\usepackage[margin=1in]{geometry}
\usepackage{amsmath,amsthm,amssymb}
\usepackage{array}
\usepackage{booktabs}
\usepackage{hyperref}
\usepackage{url}

\newcommand{\area}{\operatorname{area}}
\newcommand{\dinv}{\operatorname{dinv}}
\newcommand{\defc}{\operatorname{defc}}
\newcommand{\upmap}{\operatorname{up}}
\newcommand{\downmap}{\operatorname{down}}
\newcommand{\East}{\operatorname{East}}
\newcommand{\West}{\operatorname{West}}

\newtheorem{remark}{Remark}

\title{Computer-Assisted Well-Definedness Checks for Dyck Skeleton Strings}
\author{}
\date{}

\begin{document}
\maketitle

\section{What the computation proves}

This note explains the computer-assisted part of the skeleton-string argument
in Hawkes's 2026 paper on Dyck symmetric functions and
$q,t$-Catalan polynomials~\cite{Hawkes2026}.  The paper constructs two maps,
$\upmap$ and $\downmap$, on ordinary Dyck sequences.  In the range
\[
  \defc\leq 2n-8,
\]
these maps move one step along a fixed-deficit string:
\[
  \area(\upmap(D))=\area(D)+1,
  \qquad
  \dinv(\upmap(D))=\dinv(D)-1,
\]
with the opposite changes for $\downmap$.

The formulas defining the maps are explicit, but their well-definedness is not
formal.  A requested extraction might fail to exist; a seven-entry local move
might receive an input outside its domain; or a later reinsertion might have no
legal insertion site.  Four local lemmas in the paper rule out these failures.
The proofs combine symbolic counting with finite exhaustive computation.
Those finite checks are not merely experiments: after the symbolic argument
reduces the remaining obligations to bounded domains, the successful exhaustive
runs close those parts of the proof.

The code accompanying this note is the Appendix A and Appendix B Python from
the paper, placed in one file with a small command-line dispatcher.  Thus a
reader can reproduce the finite proof obligations without reconstructing the
appendix execution environment.

\section{Why a hybrid proof is needed}

There is a useful conceptual way to understand the proof architecture.  Let
$n'$ denote the largest size for which one could reasonably check all relevant
Dyck sequences directly, and let $n^*$ denote a threshold beyond which the
symbolic deficit-counting argument becomes easy and uniform.  These are
heuristic labels rather than constants formally defined in the paper.  The
problem is that one should expect
\[
  n'\ll n^*.
\]
A direct enumeration reaches only the genuinely small cases, while a purely
symbolic proof becomes transparent only much later.  The interval
\[
  n'<n\leq n^*
\]
is therefore too large for naive enumeration and too delicate for a short
uniform argument.

The solution is not to push either method beyond its natural range.  Instead,
the proof first derives structural restrictions on any possible failure.  A
large Dyck prefix contributes many deficit pairs, so most long configurations
are eliminated symbolically.  The few configurations that survive are encoded
by short windows, bounded prefix maxima, words with few nonzero entries, or
specific prefix forms.  Optimized programs then exhaust these reduced finite
domains.  Schematically, the proof has the form
\[
  \text{symbolic reduction}
  \; + \;
  \text{targeted exhaustive verification}
  \; = \;
  \text{complete intermediate-range argument}.
\]
The phrase ``targeted exhaustive verification'' is important: the programs do
not blindly enumerate all Dyck sequences in the difficult range.  The
mathematics determines a much smaller set of exact obligations for them to
check.

\section{Dyck sequences and the local maps}

An ordinary Dyck sequence of length $n$ is a sequence
\[
  D=(D_0,D_1,\ldots,D_{n-1})
\]
of nonnegative integers satisfying
\[
  D_0=0,
  \qquad
  D_{i+1}\leq D_i+1.
\]
Its statistics are
\[
  \area(D)=\sum_i D_i,
\]
\[
  \dinv(D)
  =\#\{(i,j):0\leq i<j<n,
      \ D_i=D_j\text{ or }D_i=D_j+1\},
\]
and
\[
  \defc(D)=\binom n2-\area(D)-\dinv(D).
\]
The paper also gives a pair-count description of $\defc$: it counts certain
``type A'' and ``type B'' pairs.  That description is the principal symbolic
tool in Appendix B.  A hypothetical failure of one of the local algorithms is
shown, whenever possible, to force more than $2n-8$ such pairs, contradicting
the assumed deficit bound.

The maps $\upmap$ and $\downmap$ are staged algorithms.  They first handle
special and skeleton cases.  Otherwise they repeatedly remove the leftmost
extractable entry, append a lowered version of each removed entry to a short
suffix, apply one of the local maps
\[
  \East_3,\ \East_5,\ \East_7
  \qquad\text{or}\qquad
  \West_3,\ \West_5,\ \West_7,
\]
and finally reinsert shifted entries.  The West maps are the reversal
conjugates of the East maps.  Three- and five-entry branches are comparatively
local.  The seven-entry branch is the difficult boundary case, and its domain
has to exclude a particular ``far-apart decomposable'' configuration.

For fixed length $n$ and deficit $d$, put
\[
  M=\binom n2,
  \qquad
  \ell=\left\lfloor\frac{M-d}{2}\right\rfloor.
\]
The intended domain of $\upmap$ consists of sequences with area at most
$\ell-1$; the intended domain of $\downmap$ consists of nonspecial sequences
with area at most $\ell$.

\section{The four local lemmas}

In the current arXiv version, the proof-critical local statements are Lemmas
5.24--5.27 of~\cite{Hawkes2026}.  Under the hypotheses
\[
  n\geq4,
  \qquad d\leq2n-8,
  \qquad \defc(D)=d,
\]
they assert the following.

\begin{center}
\begin{tabular}{>{\raggedright\arraybackslash}p{0.14\linewidth}
                >{\raggedright\arraybackslash}p{0.75\linewidth}}
\toprule
Lemma & Proof obligation \\
\midrule
5.24 & The special skeleton branches of $\upmap$ and $\downmap$ return valid
Dyck sequences; the down branch returns a full skeleton. \\
5.25 & Every extraction requested by the staged algorithms exists, and the
leftmost-extractable convention chooses it unambiguously. \\
5.26 & A seven-entry window reaching $\East_7$ or $\West_7$ is never
far-apart decomposable, so the final local map is applied inside its valid
domain. \\
5.27 & The extracted entries occur sufficiently far from the end of the
current word, and all later injection steps succeed. \\
\bottomrule
\end{tabular}
\end{center}

These lemmas feed directly into Propositions 5.29 and 5.30, which prove that
$\upmap$ and $\downmap$ are well-defined in the required lower-half domains.
Their roles are logically distinct: Lemma 5.24 handles the skeleton branches,
Lemma 5.25 guarantees the staged extraction chain, Lemma 5.26 controls the
last local window, and Lemma 5.27 guarantees that the reconstruction by
injection can actually be performed.

They also contribute to invertibility, though they are not the entire inverse
proof.  Lemma 5.27 supplies positional information showing that carried suffix
entries are not disturbed by earlier extractions.  The inverse argument also
uses the extraction--injection inverse, the local East/West inverse theorem,
and Lemma 5.23, which rules out premature termination at a smaller local
window.  Together these ingredients yield Lemma 5.32:
\[
  \downmap(\upmap(D))=D,
  \qquad
  \upmap(\downmap(D))=D
\]
whenever the corresponding maps lie in their stated domains.

The well-definedness and inverse results then give the string decomposition of
Proposition 5.34.  Repeatedly applying $\downmap$ sends every lower-half Dyck
sequence to a unique special skeleton, while repeated applications of
$\upmap$ reconstruct its string.  Theorem 5.35 converts these strings, together
with the known $q,t$-symmetry of the Catalan polynomial, into the low-deficit
skeleton formula.  Thus the appendix checks sit several steps upstream of the
final formula, but they are indispensable: without them, the maps organizing
the strings would not yet be known to exist everywhere they are used.

\section{How the finite checks divide the work}

\subsection{The residual small range}

The residual checker enumerates every Dyck sequence of lengths
\[
  4\leq n\leq7,
\]
retains the inputs satisfying the exact deficit and area hypotheses, and
follows the branch prefix of $\upmap$ or $\downmap$.  One run simultaneously
checks the small-range portions of all four local lemmas:

\begin{itemize}
\item the skeleton branches return valid words;
\item every requested extraction exists;
\item no retained input reaches a seven-window branch;
\item every extraction position used before termination satisfies the
      corresponding position bound.
\end{itemize}

A successful run prints
\begin{verbatim}
EverythingOkay = True
No East7 or West7 branch was reached for 4 <= n <= 7.
\end{verbatim}
The branch totals also match the appendix transcript: 200 eligible up calls
and 200 eligible down calls, distributed among the special, skeleton,
three-window, and five-window cases.

\subsection{The limited-nonzero check for Lemma 5.27}

The first specialized checker for Lemma 5.27 considers every Dyck sequence
with
\[
  4\leq n\leq13
\]
and at most seven nonzero entries, subject to the lemma's deficit and area
hypotheses.  This enlarged domain automatically contains every Dyck sequence
with $4\leq n\leq8$, since the initial entry is zero.  It also covers the
finite middle range left by one of the later symbolic cases.

For every eligible call, the program follows the actual $\upmap$ or
$\downmap$ branch and verifies Dyck validity, preservation of length and
deficit, the required change in area, and the extraction-position bounds.  The
recorded successful run checks 11,879 up calls and 9,486 down calls and ends
with
\begin{verbatim}
position-bound or image failures: 0
status: PASS
\end{verbatim}

\subsection{The exceptional prefix forms for Lemma 5.27}

The second Lemma 5.27 checker handles two structured prefix families that
remain after the analytic argument.  It exhausts the finite range
\[
  9\leq n\leq16.
\]
For each generated word it verifies that either the deficit bound is already
violated or the area inequality required by the putative counterexample is
violated.  It also checks the exact enumeration counts used in the proof.  A
successful run ends with
\begin{verbatim}
failures: 0
status: PASS
\end{verbatim}
This is a separate finite obligation from the limited-nonzero check; the two
programs cover different residual forms in the symbolic proof.

\subsection{The seven-window check for Lemma 5.26}

The seven-window argument is the clearest example of the hybrid method.  For
$4\leq n\leq7$, the residual checker shows that the seven-window branch is
never reached.  For $n\geq8$, Appendix B assumes that a forbidden
far-apart-decomposable window exists and then uses deficit and area inequalities
to reduce the problem.

The reduction normalizes the local seven-entry patterns, separates the cases
according to how many window entries exceed the prefix maximum, and derives
threshold tables.  Above each threshold, a descent argument rules out a
counterexample symbolically.  Below the threshold, only finitely many
integer-valued gap expansions, prefix maxima, and $(n,m)$ combinations remain.
The checker exhausts exactly those survivors.

The run constructs 7,194 East windows and 7,194 West windows.  It verifies the
expected threshold tables, the structural lower bound
$\operatorname{id}_{\mathrm{mid}}\geq10$, and the four finite Case/Side
searches.  The expected child and triple counts are
\[
\begin{array}{c|r|r}
&\text{children}&\text{triples}\\
\hline
\text{Case 1 East}&2473&9919\\
\text{Case 1 West}&2911&10311\\
\text{Case 2 East}&3860&715\\
\text{Case 2 West}&4827&1756
\end{array}
\]
and every search reports zero problems.  The final line is
\begin{verbatim}
SUCCESS: East7/West7 seven-window verification passed.
\end{verbatim}
The program is therefore not enumerating all length-$n$ Dyck sequences.  It is
checking the bounded local configurations that survive a substantial symbolic
reduction.

\section{Using \texttt{code.py}}

The file \texttt{code.py} embeds the Appendix A core routines and the four
Appendix B checker listings from the paper.  The only repository-specific
addition is a dispatcher that executes a selected listing in a fresh shared
namespace.  The available commands are
\begin{verbatim}
python code.py residual
python code.py limited-nonzero
python code.py prefix
python code.py seven-window
python code.py all
\end{verbatim}
The default is \texttt{all}.  On the machine used to prepare this note, the
four checks together complete in a few seconds; runtimes vary by machine.

The embedded paper listings can also be written back out as separate files:
\begin{verbatim}
python code.py all --extract appendix_listings
\end{verbatim}
This option is only a convenience for inspection.  It does not alter the code
that is executed.

The checkers require only the Python standard library.  They must be run with
ordinary assertion checking enabled.  In particular, do not use
\begin{verbatim}
python -O code.py all
\end{verbatim}
because several assertions are part of the verification.  The dispatcher
rejects optimized mode for this reason.

\section{What the computation does not do}

The successful runs do not independently prove the skeleton-string theorem.
They assume the definitions, reductions, threshold inequalities, and domain
translations established in the written argument.  Conversely, they should
not be described as optional numerical evidence: on the bounded domains left
by those reductions, the exhaustive checks are proof steps.

The correct division of labor is therefore:
\begin{itemize}
\item the paper's symbolic reasoning proves that every possible failure lies
      either in an analytically impossible family or in one of the encoded
      finite domains;
\item the code exhausts those finite domains and finds no failure;
\item the local lemmas then establish well-definedness and supply part of the
      control needed for invertibility;
\item the well-defined inverse maps produce the lower-half string
      decomposition used in the final $q,t$-Catalan formula.
\end{itemize}
If the definitions or symbolic reductions in the paper are changed, the
correspondence between those arguments and the finite checker domains must be
reviewed again.

\begin{thebibliography}{9}

\bibitem{Hawkes2026}
Graham Hawkes,
\emph{Dyck Symmetric Functions and Applications to $q,t$-Catalan Polynomials},
arXiv:2605.13003, 2026.
\url{https://arxiv.org/abs/2605.13003}

\end{thebibliography}

\end{document}
